\section{KHÓ KHĂN VÀ CÁCH XỬ LÝ}\label{sec:kho_khan}

Bảng~\ref{tab:kho_khan} tổng hợp các vấn đề gặp trong tuần. Mỗi vấn đề được chẩn đoán theo lớp và có bằng chứng (log, kết quả đo) trước khi sửa.

\begingroup
\small
\setlength{\tabcolsep}{4pt}
\begin{longtable}{|P{3.4cm}|P{4.6cm}|P{4.8cm}|}
    \caption{Khó khăn trong tuần và cách xử lý}\label{tab:kho_khan} \\
    \hline
    \textbf{Hiện tượng} & \textbf{Nguyên nhân} & \textbf{Cách xử lý / trạng thái} \\
    \hline
    \endfirsthead
    \multicolumn{3}{c}{\textit{Bảng~\ref{tab:kho_khan} (tiếp theo)}} \\
    \hline
    \textbf{Hiện tượng} & \textbf{Nguyên nhân} & \textbf{Cách xử lý / trạng thái} \\
    \hline
    \endhead
    \hline
    \endfoot
    Controller Ryu khựng, nhiều switch báo mất kết nối điều khiển & Máy chủ EVE-NG dùng chung tải rất cao khi chạy cùng lúc cụm SDN và cụm SD-WAN & Tắt mềm cụm SD-WAN khi đo SDN: tải máy chủ giảm mạnh, hết hiện tượng kết nối lại, tỉ lệ kiểm tra sức khỏe không đạt giảm từ 23\% xuống 7--10\% \\
    \hline
    Controller vẫn khựng 2--3 lần mỗi phút, mỗi lần 0,6--1,5~s & Tiến trình Ryu không bị treo; CPU ảo của máy ảo controller được đánh thức muộn do máy chủ dùng chung & Ghim CPU cho máy ảo controller (giảm nhẹ); đưa phần chuyển đường quan trọng xuống data plane bằng fast-failover. Tinh chỉnh toàn máy chủ cần xin phép trước \\
    \hline
    Mất một uplink Access gián đoạn 6--8~s & Thời gian ân hạn 3,5~s cho mỗi lần controller khựng đẩy việc kết luận liên kết chết lùi lại & BFD + group OpenFlow fast-failover; còn 1--2 gói \\
    \hline
    Bộ ping trong controller báo mất gói khi data plane vẫn thông & Bộ ping chạy cùng tiến trình controller nên bị ảnh hưởng bởi các lần khựng & Đo phục hồi bằng ping từ máy trạm VPC trong khi cắt tap thật \\
    \hline
    Sai số khi đo hiệu năng & Đích mặc định là một máy trạm đang tắt & Đổi đích sang máy chủ ở Server Farm \\
    \hline
    Biểu đồ trên Console co dần sau mỗi lần cập nhật & Hàm vẽ đọc lại chiều cao canvas đã nhân tỉ lệ điểm ảnh & Tách kích thước hiển thị khỏi kích thước canvas; kiểm tra ở nhiều tỉ lệ \\
    \hline
    API ping của vManage trả HTTP 400 ``Invalid data length value'' & vManage 20.10 chỉ chấp nhận một số tổ hợp tham số kích thước gói & Thử lần lượt các dạng tham số và nhớ dạng chạy được \\
    \hline
    Các switch OVS khởi động rất chậm (khoảng 10 phút) & Tải máy chủ cao khi bật lại lab & Bật theo thứ tự đã kiểm chứng: dịch vụ nền, controller, Distribution, Access rồi mới tới máy trạm \\
\end{longtable}
\endgroup
